\chapter{迁移学习}
\label{chap:transfer}
\begin{quote}\small
\textit{本章摘要。}迁移学习利用已有数据、表示或参数降低目标任务的学习成本。本章先区分任务变化与分布偏移，再从风险重加权、表示对齐和参数适配推导主要算法，最后讨论测试时适配与负迁移。选择依据是哪些假设仍成立，而不是预训练模型有多大。
\end{quote}

上一章把大量历史观测压缩成一个可复用的编码器，但“可复用”只是潜力，不是新设备上的性能承诺。旧相机学到的边缘结构可能仍然有用，照明统计可能需要更新，而不同产品的缺陷定义甚至可能不再相同。本章因此从参数初始化走向目标风险：不是询问源模型有多强，而是询问源数据中的哪些关系在目标环境中仍成立。

这里存在三种不同的干预。若预测关系不变而采样频率改变，可重加权样本；若不同观测方式背后存在共享语义，可尝试对齐表示；若目标任务需要新的决策规则，则必须更新分类头或部分骨干。三种干预对应不同假设，不能把它们看成可互换的技巧。本章先推导风险如何从源分布转换到目标分布，再解释统计对齐与域对抗到底消除了什么差异，随后讨论有限标签与计算预算下的参数适配，并用负迁移说明这些方法的边界。

假设两台相机的图像颜色变得更接近了，还要检查裂纹和正常纹理是否仍能区分。对齐指标下降，只说明某种分布差异缩小，不能直接说明目标设备上的故障更容易识别。章末再考虑连续更新：用新设备数据改好模型后，旧设备上的故障还能不能识别？下一章的持续学习就从这里开始。

\section{源域知识与目标域风险}
\begin{figure}[htbp]
\centering
\begin{tikzpicture}[node distance=0.75cm and 0.7cm]
 \node[idi input, minimum width=2.0cm] (source) {源域\\$P_s(X,Y)$};
 \node[idi process, right=of source, minimum width=2.0cm] (repr) {共享表示\\$F_\theta(X)$};
 \node[idi output, right=of repr, minimum width=2.0cm] (target) {目标域\\$P_t(X,Y)$};
 \node[idi state, below=0.8cm of repr, minimum width=2.3cm] (adapt) {重加权、对齐\\或参数适配};
 \draw[idi arrow] (source)--(repr); \draw[idi arrow] (repr)--(target);
 \draw[idi arrow] (source)--(adapt); \draw[idi arrow] (target)--(adapt);
 \draw[idi feedback] (adapt)--(repr);
\end{tikzpicture}
\caption{迁移学习的核心是把源域知识送入目标域，同时检验哪些分布或任务假设仍然成立。}
\label{fig:transfer-source-target}
\end{figure}
图\ref{fig:transfer-source-target}将源域知识、适配过程和目标域评价分开；共享参数并不意味着分布相同，目标域验证必须检验迁移后是否降低了实际风险。
设源域与目标域联合分布分别为$P_s(X,Y)$与$P_t(X,Y)$，目标是降低$R_t(f)=\E_{P_t}\ell(f(X),Y)$。迁移学习允许域或任务发生变化，领域自适应通常保持预测任务并处理域差异，领域泛化则要求在训练时未见的域上工作\cite{pan2010transfer}。新相机质检可保持缺陷类别不变，跨产品迁移却可能改变标签空间；二者不能使用同一适配规则。

协变量偏移假定$P_s(Y\mid X)=P_t(Y\mid X)$而输入边缘分布变化；标签偏移假定$P_s(X\mid Y)=P_t(X\mid Y)$而类别比例变化；概念变化则涉及预测关系本身。实际数据可能同时存在多种变化。仅观察无标签目标输入通常无法确定标签机制是否改变，因此无监督适配不能替代目标标签审核。

\section{实例重加权与分布对齐}
\paragraph{重要性加权。}
在协变量偏移且目标分布支撑被源分布覆盖时，目标风险可写为
\begin{equation}
 R_t(f)=\E_{P_s}\left[w(X)\ell(f(X),Y)\right],\qquad
 w(x)=\frac{p_t(x)}{p_s(x)}.
\end{equation}
先用源、目标输入估计密度比，再最小化加权源域损失。核均值匹配直接使加权源域特征均值接近目标均值，避免分别估计高维密度\cite{huang2007kmm}。若用域分类器估计比值，必须校正训练时源、目标样本的先验比例。大权重会使少数样本主导梯度，截断权重可以降低方差但引入偏差。有效样本量$(\sum_iw_i)^2/\sum_iw_i^2$可用于判断加权是否退化。

\paragraph{风险恒等式、密度比估计与方差。}
以离散标签和连续输入为例，协变量偏移要求对目标几乎处处的$x$，两域条件分布相同，且$p_t(x)>0$时$p_s(x)>0$。于是
\begin{align}
 R_t(f)&=\int\sum_y\ell(f(x),y)p_t(y\mid x)p_t(x)\,\mathrm dx\\
 &=\int\sum_y\ell(f(x),y)p_s(y\mid x)p_s(x)
 \frac{p_t(x)}{p_s(x)}\,\mathrm dx.\label{eq:transfer-source-risk-integral}
\end{align}
式\eqref{eq:transfer-source-risk-integral}才是可从源标签估计的量。支撑覆盖不是技术性装饰：源域从未出现的输入区域没有标签项可重用，任何有限权重都不能补上这部分风险。

令域指示$D=1$表示目标，域分类器训练混合中的先验为$\pi_t$与$\pi_s=1-\pi_t$，其校准输出为$q(x)=P(D=1\mid x)$。由贝叶斯公式
\begin{equation}
 \frac{q(x)}{1-q(x)}=\frac{\pi_t p_t(x)}{\pi_s p_s(x)},\qquad
 w(x)=\frac{\pi_s}{\pi_t}\frac{q(x)}{1-q(x)}.
\end{equation}
若训练时人为平衡两域，先验就是该平衡采样的比例，而不是原始数据库大小的比例。$q$接近$1$使估计非常不稳定；高域分类准确率此时可能表明重叠不足，而不是重加权更可靠。

对独立源样本，使用真实权重的估计$\widehat R=n_s^{-1}\sum_iw_i\ell_i$在可积时无偏；有限方差时其方差为$\Var_s(w\ell)/n_s$。若截断为$\bar w=\min(w,c)$，且$0\leq\ell\leq L_{\max}$，则
\begin{equation}
 0\leq R_t-\E_s[\bar w\ell]
 =\E_s[(w-c)_+\ell]\leq L_{\max}\E_s[(w-c)_+].
\end{equation}
这显式展示了截断的偏差来源。自归一化估计$\sum_iw_i\ell_i/\sum_iw_i$也通常具有有限样本偏差，但可改善数值稳定性。有效样本量来自归一化权重$\alpha_i=w_i/\sum_jw_j$的$1/\sum_i\alpha_i^2$：全部相等时为$n_s$，一个样本几乎占全重时接近$1$；它是权重集中的诊断量，不是所有损失下精确的独立样本数。

标签偏移下，权重改为类别先验比$p_t(y)/p_s(y)$。这个结论来自$p_t(x,y)=p_s(x\mid y)p_t(y)=p_s(x,y)p_t(y)/p_s(y)$，要求目标出现的类别在源域具有正概率。若希望从无标签目标数据估计先验，可先固定一个分类器，定义源验证集上的混淆矩阵$C_{ky}=P_s(\widehat Y=k\mid Y=y)$。标签偏移意味着该矩阵在目标域不变，目标预测类别频率$\mu_k=P_t(\widehat Y=k)$满足$\mu=C\pi_t$，其中$\pi_t\in\R^K$是目标类别先验。可在$\pi_t\geq0$、$\mathbf1^\top\pi_t=1$约束下拟合这个线性关系。矩阵列若近似相同，分类器不能区分类别，先验反演便病态；若$P(X\mid Y)$也改变，则关系本身失效。它不能直接与协变量偏移权重互换。若目标域出现源域从未包含的故障，仅靠重加权不能产生缺失的类别证据。

\paragraph{Maximum Mean Discrepancy (MMD) 与 Correlation Alignment (CORAL)。}
深度适配网络在特征$z=F_\theta(x)$上同时优化源域监督损失和域差异\cite{long2015dan}：
\begin{equation}
 \mathcal L=\mathcal L_s+\lambda\left\|\frac1{n_s}\sum_i\varphi(z_i^s)
 -\frac1{n_t}\sum_j\varphi(z_j^t)\right\|_{\mathcal H}^2.\label{eq:transfer-mmd-objective}
\end{equation}
核技巧将平方范数展开为源源、目标目标和源目标三组核值，无需显式构造$\varphi$。训练时从两域抽取批次，计算分类与对齐损失，再共同更新编码器。核带宽决定比较的尺度，多核 MMD 组合不同带宽以减少对单一尺度的依赖。

\paragraph{核均值差的实际计算。}
设$k(u,v)=\langle\varphi(u),\varphi(v)\rangle_{\mathcal H}$为正定核，$z\in\R^d$。将式\eqref{eq:transfer-mmd-objective}中的平方范数按$\|a-b\|^2=\langle a,a\rangle+\langle b,b\rangle-2\langle a,b\rangle$展开，得到
\begin{align}
 \widehat{\mathrm{MMD}}_b^2
 &=\frac1{n_s^2}\sum_{i,i'}k(z_i^s,z_{i'}^s)
 +\frac1{n_t^2}\sum_{j,j'}k(z_j^t,z_{j'}^t)\\
 &\quad-\frac2{n_sn_t}\sum_{i,j}k(z_i^s,z_j^t).
\end{align}
这是含对角项的非负、有偏估计。若两域样本分别独立同分布且$n_s,n_t>1$，去掉前两项中相同下标的配对，并将分母换为$n_s(n_s-1)$和$n_t(n_t-1)$，可得到总体平方 MMD 的无偏估计；后者有限样本下可能为负，不能误认为程序必定出错。采用特征核时，总体 MMD 为零可以识别分布相同；线性核只检查均值相同。无论采用哪种核，比较对象都是$P_s(Z)$与$P_t(Z)$，并没有自动检查$P_s(Y\mid Z)$与$P_t(Y\mid Z)$。

CORAL 对齐二阶统计量\cite{sun2016coral}：
\begin{equation}
 \mathcal L_{\mathrm{CORAL}}=\frac{1}{4d^2}\|C_s-C_t\|_F^2,
\end{equation}
其中$C_s,C_t$为$d$维表示的批次协方差。具体地，把样本按行写成$Z_s\in\R^{n_s\times d}$，令$H_s=I_{n_s}-\mathbf1\mathbf1^\top/n_s$，则$C_s=Z_s^\top H_sZ_s/(n_s-1)$，目标域同理。由于$C_s-C_t$对称，矩阵微分给出
\begin{equation}
 \nabla_{Z_s}\mathcal L_{\mathrm{CORAL}}
 =\frac{H_sZ_s(C_s-C_t)}{d^2(n_s-1)}.
\end{equation}
可见梯度来自中心化表示与协方差差异的乘积，源、目标均值本身不在该损失中；如果还需均值对齐，应另行添加对应项。$n_s=1$时样本协方差未定义，小批次下秩至多为$n_s-1$，也解释了协方差估计为何依赖批次组织。它计算直接，但小批次协方差不稳定，而且二阶一致不代表类别条件分布一致。MMD 与 CORAL 都可能把不同类别拉近；存在目标标签时应检查类条件结构，使用伪标签时则需控制错误对齐的累积。

\paragraph{域对抗适配。}
Domain-Adversarial Neural Network (DANN) 用域判别器$D$识别表示来自哪一个域\cite{ganin2016dann}。若域损失为交叉熵，目标写为
\begin{equation}
 \min_{F,C}\max_D\;\mathcal L_s(C(F(x_s)),y_s)
 -\lambda\mathcal L_d(D(F(x)),d).
\end{equation}
域判别器最小化自己的交叉熵，编码器通过梯度反转最大化该损失，同时维持源域分类能力。实现时不能将两个模块都沿同一符号更新。域准确率接近随机只是对齐的诊断信号，并非目标分类正确的证据；特征全部坍缩也可使域不可区分。

\paragraph{域不可区分为何不是语义正确。}
为了看清目标的符号，采用等域权重的域损失
\begin{equation}
 \mathcal L_d=-\E_s\log(1-D(Z))-\E_t\log D(Z).
\end{equation}
固定编码器，在具有特征密度的情况下，对每个$z$逐点最小化，得到$D^*(z)=p_t(z)/(p_s(z)+p_t(z))$。代回可得
\begin{equation}
 \min_D\mathcal L_d=\log4-2\operatorname{JS}(P_s^Z\|P_t^Z),
\end{equation}
其中$\operatorname{JS}(P\|Q)=\tfrac12\KL(P\|M)+\tfrac12\KL(Q\|M)$，$M=(P+Q)/2$。因此在判别器充分优化这一理想条件下，编码器最大化域损失相当于减小特征边缘分布差异。若实现中再把两域损失除以$2$，常数和系数也相应减半，方向不变。

但令两域特征都均匀取$-1,+1$，源域标签满足$Y=\mathbf1[Z=1]$，目标域却满足$Y=\mathbf1[Z=-1]$。两域特征分布完全相同，任何域判别器都只能猜测，而源分类器在目标域全部判错。这个反例说明真正需要的是可共享的预测关系，而不只是看不出设备来源。二分类的零一风险可用理论形式
\begin{equation}
 R_t(h)\leq R_s(h)+\tfrac12d_{\mathcal H\Delta\mathcal H}(P_s^X,P_t^X)+\lambda^*,
 \qquad \lambda^*=\min_{h\in\mathcal H}(R_s(h)+R_t(h))
\end{equation}
理解这一点：差异量定义为$2\sup_{h,h'\in\mathcal H}|P_s(h\ne h')-P_t(h\ne h')|$，检查假设对在两域的分歧概率；$\lambda^*$则表达同一分类器同时做好两个域的可能性。这是总体风险界，有限样本估计还需要复杂度误差项；域对抗损失也不等于这个差异量的精确值。对齐可以降低某种域差异，却不能从无标签输入中证明$\lambda^*$很小。

条件对齐将类别预测纳入域判别输入，以区分同一域中的不同语义。多源适配对不同源域分别估计可信度，避免样本最多的源主导目标；部分域适配需降低源域独有类别的权重，开放集适配则必须识别目标域未知类别。这些变体改变的是标签空间假设，不能只通过调大域对齐系数实现。

\section{预训练、微调与参数高效适配}
第\ref{chap:pretraining}章讨论如何从预训练数据构造可复用表示，本节转向已有表示如何适应目标任务。
预训练提供初始化，但目标任务仍决定哪些层需要改变。线性探测冻结编码器只训练头，可检验原表示是否已具有判别性；逐层解冻先调整高层，再判断低层滤波是否需要变化；全量微调自由度最大，也更容易在少量目标样本上过拟合。比较时应固定目标标签量、预训练来源与超参数搜索预算。

Adapter 在冻结骨干中插入低维瓶颈。对输入表示$h$，一个残差适配器可写为$h'=h+W_{up}\sigma(W_{down}h)$，只训练新增参数\cite{houlsby2019adapter}。LoRA 则令权重更新$\Delta W=BA$，以秩$r$限制更新自由度\cite{hu2021lora}。二者都降低每个目标域的存储开销，但低秩不意味着任何域变化都能被表达。提示学习优化输入或中间层的提示向量，适合保持基础模型不变的条件化使用；其可迁移性仍取决于基础模型是否理解目标语义。

\paragraph{低秩更新的维度、梯度与容量。}
令冻结的线性层为$W_0\in\R^{d_{out}\times d_{in}}$，输入$x\in\R^{d_{in}}$。LoRA 使用$A\in\R^{r\times d_{in}}$、$B\in\R^{d_{out}\times r}$及尺度$s=\alpha/r$，输出为
\begin{equation}
 y=(W_0+sBA)x.
\end{equation}
可训练参数由$d_{out}d_{in}$降为$r(d_{out}+d_{in})$，更新秩不超过$r$。若$G=\partial\mathcal L/\partial W_{eff}$，矩阵链式法则给出
\begin{equation}
 \nabla_B\mathcal L=sGA^\top,\qquad
 \nabla_A\mathcal L=sB^\top G.
\end{equation}
这解释了常见的随机初始化$A$、零初始化$B$：初始$BA=0$保证原输出不变，同时$B$可以收到非零梯度；若两者都初始化为零，则两条梯度都为零，低秩分支无法启动。假如理想全量更新$\Delta W^*$的奇异值为$\sigma_1\geq\sigma_2\geq\cdots$，最佳秩$r$矩阵近似的平方 Frobenius 误差为$\sum_{j>r}\sigma_j^2$。因此低秩适配隐含的是有用更新可被少数方向近似，而非原权重本身低秩；该矩阵逼近结论也不等于网络任务风险的误差界。

对 Adapter，若$h\in\R^d$、瓶颈宽度为$r$，则$W_{down}\in\R^{r\times d}$、$W_{up}\in\R^{d\times r}$，忽略偏置时有$2dr$个参数。它的非线性瓶颈通常不能像单个线性 LoRA 更新那样直接合并进原权重。选择应同时看可表达变化、路由和推理限制，而不只看参数比率。

参数量、训练显存与推理开销应分别报告。冻结权重仍可能需要存储用于反向传播的激活；多个 Adapter 需要路由规则，LoRA 合并则需要确认目标域选择是否明确。不能把参数高效方法自动视为持续学习算法：若所有域共享并反复覆盖同一个低秩更新，仍会遗忘旧域。

\section{测试时适配与负迁移}
测试时适配只使用部署输入调整有限参数。Tent 最小化预测熵，并在其设定下更新归一化层仿射参数\cite{wang2021tent}：
\begin{equation}
 \min_{\eta}\frac1{|B|}\sum_{x\in B}
 \left[-\sum_kp_\eta(k\mid x)\log p_\eta(k\mid x)\right].
\end{equation}
先保留训练好的模型，按部署批次计算熵，再对允许更新的参数执行梯度步。低熵可能是正确判断，也可能是错误自信。类别极不平衡、批次很小或流中含未知故障时，熵最小化可能将预测推向单一类别。逐批重置与持续累积更新是不同规则，必须明确是否保留适配状态。

\paragraph{熵下降为什么会强化错误。}
令类别 logits 为$a\in\R^K$、$p=\softmax(a)$。由$\partial p_k/\partial a_j=p_k(\mathbf1[k=j]-p_j)$，可逐项得到
\begin{equation}
 \frac{\partial H(p)}{\partial a_j}=-p_j(\log p_j+H(p)).
\end{equation}
二分类时令$p=\sigma(a)$，则$\mathrm dH/\mathrm da=p(1-p)\log((1-p)/p)$。若$p>1/2$，梯度下降会继续增大$a$和$p$；若$p<1/2$则进一步压低$p$。没有标签参与时，这个更新没有能力识别当前偏向是否正确，只会强化已有偏向。熵适配因而依赖初始预测大体可信、低密度区域适合作决策边界等额外条件。应限制步数和漂移幅度，保留重置机制，并使用独立监测数据检验校准与少数类召回；这些操作降低风险，但仍不能构成无标签正确性保证。

负迁移指相对于合适的无迁移基线，使用源知识反而提高目标风险。应比较从零训练、线性探测、全量微调和轻量适配，并绘制随目标标签量变化的学习曲线。领域自适应理论中的分布差异项之外，还存在源、目标任务能否被同一假设同时解释的误差项\cite{bendavid2010theory}；只减小域差异并不足够。

换了新相机后，先查清楚是亮度、镜头和噪声变了，还是“什么算缺陷”的定义变了。前一类问题可以先试归一化适配和线性探测；后一类问题需要重新核对标签。若规则允许，可以用新相机的无标签图像做适配，但测试标签只能留到最终评价。报告各类缺陷召回、未知输入拒识、校准、训练成本和发生负迁移的设备组。若还要长期维护旧相机能力，就需要下一章的持续学习机制。

\paragraph{本章小结。}
迁移学习依次回答源知识是否相关、分布差异能否被估计、哪些参数应更新。重加权改变样本贡献，成立前提是正确的偏移假设和支撑覆盖；分布对齐改变表示，却必须额外保护类别语义；微调和低秩适配改变预测器，其容量与标签预算应匹配。测试时熵下降进一步说明，没有标签的自洽改进可能只是错误自信。所有路线都需要目标域证据与无迁移基线，不能以对齐损失下降代替目标风险下降。

实际选方法时，先核对新旧数据的标签含义是否相同，再看目标设备有多少标签、能否使用它的无标签记录，最后选择足够解决问题的适配办法。如果系统仍需服务旧相机，更新后就要一起复测；新相机得分提高、旧相机却明显退步，不能算整个更新任务已经完成。下一章据此讨论持续学习：数据不断到来时，怎样学习新情况、保留旧能力，并控制存储和计算开销。